% ═══════════════════════════════════════════════════════════════
% Section 5: NISS — Neural Impact Scoring System
% ═══════════════════════════════════════════════════════════════

\section{NISS: Neural Impact Scoring System}
\label{sec:niss}

The Neural Impact Scoring System (NISS) is a proposed CVSS v4.0 extension
designed to conform with FIRST.org's official extension mechanism (User Guide
\S3.11)~\cite{first2023cvss4userguide}. NISS adds six metrics that capture
dimensions CVSS was never designed to express: biological tissue damage, cognitive
reconnaissance, cognitive disruption, consent boundary violations, damage
reversibility, and neuroplastic consequences.

NISS scores impact only; it has no likelihood or exploitability term and is meant to sit
beside a CVSS base score, not to replace it. It is also the technique-level input to a
proposed multi-layer Neurosecurity Score (NSv2.1), which is unvalidated and outside the
scope of this paper.

\subsection{Gap Analysis: Why CVSS Alone Is Insufficient}

Each TARA technique carries a CVSS v4.0 base vector and a gap group recording how much of
its impact CVSS fails to capture. Table~\ref{tab:gap} gives the counts published in v1.5
beside those recomputed from the registry file committed with v1.5.

\begin{table}[H]
\centering
\caption{CVSS v4.0 gap groups across the \SnapshotTechniques{} snapshot techniques: as
published in v1.5, and recomputed from the committed snapshot file.}
\label{tab:gap}
\small
\begin{tabularx}{\textwidth}{@{}c X r r@{}}
\toprule
\textbf{Group} & \textbf{Gap Description} & \textbf{v1.5 text} & \textbf{Snapshot file} \\
\midrule
1 & CVSS captures most impact; NISS adds nuance        & 13 & 14 \\
2 & CVSS captures exploitability but misses part of the impact & 31 & 30 \\
3 & CVSS cannot express the primary impact             & 64 & 59 \\
--- & No gap group assigned                            & 1  & 6 \\
\midrule
  & \textbf{Groups 1--3} & \textbf{\SnapshotGapTechniques{} (\SnapshotGapPercent\%)} & \textbf{103 (94.5\%)} \\
\bottomrule
\end{tabularx}
\end{table}

The v1.5 counts could not be reproduced from the registry file committed with that
version, which gives 59 techniques (54.1\%) in Group~3 and 103 (94.5\%) in any group; this
revision uses the reproducible counts in its text and retains the published ones in the
table for the record. Group~1 techniques are mostly captured by CVSS, so the share that
\emph{needs} NISS is better read as Groups~2 and~3, 89 techniques (81.7\%). In Group~3 the
most severe consequence is tissue damage, disruption of neural function or lasting
harm---dimensions for which CVSS has no metric. Live: \LiveGapTechniques{} of
\LiveTechniques{} (\LiveGapPercent\%) carry a gap group.

\subsection{Extension Metrics}

NISS defines six extension metrics, each with a graduated value set. The
metrics are designed to be orthogonal to CVSS base metrics: they capture impact
dimensions that exist only because the target system interfaces with biological
neural tissue.

\subsubsection{BI: Biological Impact}

Direct harm to neural tissue, organs, or physiological function. This dimension
has no equivalent in CVSS.

\begin{table}[H]
\centering
\small
\begin{tabular}{@{}l l l p{7cm}@{}}
\toprule
\textbf{Value} & \textbf{Label} & \textbf{Score} & \textbf{Description} \\
\midrule
N & None     & 0.0  & No tissue interaction or physical harm \\
L & Low      & 3.3  & Temporary discomfort, minor sensory disruption, reversible tissue stress \\
H & High     & 6.7  & Significant tissue damage, seizure induction, involuntary motor activation \\
C & Critical & 10.0 & Life-threatening or permanently disabling neural harm \\
\bottomrule
\end{tabular}
\end{table}

\subsubsection{CR: Cognitive Reconnaissance}

Degree of unauthorised read access: decoding of speech or motor intent by models trained
on a cooperating user, inference from neural data, intent extraction. CVSS treats such
reads only as loss of data confidentiality. CR maps to the Cognitive Liberty and Mental
Privacy neurorights dimensions.

\begin{table}[H]
\centering
\small
\begin{tabular}{@{}l l l p{7cm}@{}}
\toprule
\textbf{Value} & \textbf{Label} & \textbf{Score} & \textbf{Description} \\
\midrule
N & None     & 0.0  & No cognitive read access \\
L & Low      & 3.3  & Partial intent inference, coarse emotional state detection \\
H & High     & 6.7  & Decoding of intent, identity inference, or detailed neural data extraction \\
C & Critical & 10.0 & Sustained, broad inference from neural data in real time \\
\bottomrule
\end{tabular}
\end{table}

\subsubsection{CD: Cognitive/Functional Disruption}

Degree of unauthorized disruption to cognitive processing, sensory perception,
motor output, or autonomic regulation---including perception manipulation,
identity modification, cognitive coercion, and false stimuli injection. CD maps
to the Mental Integrity neurorights dimension and extends to ICD-10-CM chapters
G~(nervous system), H~(eye/ear), and R~(symptoms/signs).

\begin{table}[H]
\centering
\small
\begin{tabular}{@{}l l l p{7cm}@{}}
\toprule
\textbf{Value} & \textbf{Label} & \textbf{Score} & \textbf{Description} \\
\midrule
N & None     & 0.0  & No cognitive write access or alteration \\
L & Low      & 3.3  & Minor perceptual distortion, detectable interference \\
H & High     & 6.7  & Perception manipulation, involuntary behavioral changes \\
C & Critical & 10.0 & Cognitive coercion, identity modification, or complete loss of cognitive autonomy \\
\bottomrule
\end{tabular}
\end{table}

\subsubsection{CV: Consent Violation}

Degree of violation of informed consent or cognitive autonomy. Ordered by
severity: covert (implicit) violations are worse than detectable (explicit) ones.

\begin{table}[H]
\centering
\small
\begin{tabular}{@{}l l l p{7cm}@{}}
\toprule
\textbf{Value} & \textbf{Label} & \textbf{Score} & \textbf{Description} \\
\midrule
N & None              & 0.0  & Operating within explicitly consented boundaries \\
P & Partial           & 3.3  & Action exceeds scope but subject retains some awareness \\
E & Explicit          & 6.7  & Direct violation of consent boundaries, but detectable \\
I & Implicit (covert) & 10.0 & Covert manipulation the patient cannot detect or refuse \\
\bottomrule
\end{tabular}
\end{table}

\subsubsection{RV: Reversibility}

Whether the damage can be undone. IT assets can be restored from backup. Neural
tissue cannot be rebooted.

\begin{table}[H]
\centering
\small
\begin{tabular}{@{}l l l p{7cm}@{}}
\toprule
\textbf{Value} & \textbf{Label} & \textbf{Score} & \textbf{Description} \\
\midrule
F & Full        & 0.0  & Effects fully reverse when attack stops \\
T & Temporary   & 3.3  & Effects reverse over hours to days \\
P & Partial     & 6.7  & Some effects permanent, some reversible \\
I & Irreversible & 10.0 & Permanent neural tissue destruction or cognitive change \\
\bottomrule
\end{tabular}
\end{table}

\subsubsection{NP: Neuroplasticity}

Whether the attack exploits or induces neuroplastic changes---the brain's ability
to rewire itself. This has no digital equivalent.

\begin{table}[H]
\centering
\small
\begin{tabular}{@{}l l l p{7cm}@{}}
\toprule
\textbf{Value} & \textbf{Label} & \textbf{Score} & \textbf{Description} \\
\midrule
N & None       & 0.0  & No neuroplastic effect \\
T & Temporary  & 3.3  & Short-term synaptic changes that decay within hours--weeks \\
P & Partial    & 6.7  & Moderate synaptic reorganization persisting weeks to months \\
S & Structural & 10.0 & Long-term or permanent neural pathway changes \\
\bottomrule
\end{tabular}
\end{table}

\subsection{PINS Flag}

NISS introduces the Potential Impact to Neural Safety (PINS) flag---a
binary indicator triggered when:

\begin{equation}
\text{PINS} = \begin{cases}
  \texttt{true} & \text{if } \text{BI} \geq \text{High} \;\lor\; \text{RV} = \text{Irreversible} \\
  \texttt{false} & \text{otherwise}
\end{cases}
\end{equation}

A PINS flag calls for safety review regardless of overall score. In the snapshot, 33 of
\SnapshotTechniques{} techniques are PINS-flagged (30.3\%); in the live registry,
\LivePinsFlagged{} of the \LiveNissScored{} scored techniques are.

\subsection{Scoring Formula}

The NISS score is computed as the weighted mean of the six metric scores, rounded up to
one decimal place:

\begin{equation}
\text{NISS} = \frac{w_{\text{BI}} \cdot \text{BI} + w_{\text{CR}} \cdot \text{CR} + w_{\text{CD}} \cdot \text{CD} + w_{\text{CV}} \cdot \text{CV} + w_{\text{RV}} \cdot \text{RV} + w_{\text{NP}} \cdot \text{NP}}
              {w_{\text{BI}} + w_{\text{CR}} + w_{\text{CD}} + w_{\text{CV}} + w_{\text{RV}} + w_{\text{NP}}}
\label{eq:niss}
\end{equation}

In the default profile, the weights are $w_{\text{BI}} = w_{\text{CV}} =
w_{\text{RV}} = w_{\text{NP}} = 1.0$ and $w_{\text{CR}} = w_{\text{CD}} = 0.5$.
This normalization ensures that the split of the former Cognitive Integrity (CG)
metric into Cognitive Reconnaissance (CR) and Cognitive Disruption (CD) in v1.1
does not inflate the cognitive dimension's influence: at $0.5 + 0.5 = 1.0$ out of
a total weight sum of~5.0, cognition retains its original 20\% share from the
five-metric v1.0 formula. The CR/CD split reflects the fundamental distinction
between read attacks (inference from neural data) and write attacks
(perception manipulation, cognitive coercion)---attack vectors with different
mitigations that map to different neurorights dimensions.
NISS supports four context profiles with differential weights that override
the default normalization:

\begin{itemize}
  \item \textbf{Clinical}: BI, CD and RV at 2.0, others at 1.0
  \item \textbf{Research}: CR and CV at 2.0, CD and NP at 1.5
  \item \textbf{Consumer}: CR and CV at 2.0, others at 1.0
  \item \textbf{Military}: BI, CR and CD at 2.0, RV and NP at 1.5, CV at 0.5
\end{itemize}

\noindent All stored scores use the default weights. The military profile halves the
weight of consent, a choice with ethical content that has not been reviewed.

\subsection{Vector Format}

The NISS vector rides alongside the CVSS v4.0 base vector:

\begin{lstlisting}
CVSS:4.0/AV:N/AC:L/AT:N/PR:N/UI:N/
  VC:H/VI:H/VA:H/SC:N/SI:N/SA:N
NISS:1.1/BI:H/CR:H/CD:C/CV:I/RV:P/NP:S
\end{lstlisting}

This dual-vector architecture means security teams can triage using familiar CVSS
scores while BCI-specific teams see the neural dimensions that determine whether a
vulnerability is a software bug or a patient safety emergency.

\subsection{Scoring States and Stated Limits}
\label{sec:niss-limits}

A NISS score is six judgements about impact. Scoring a technique whose sources support only
one of them means inventing the other five, so the registry now distinguishes three
\emph{scoring states}:

\begin{itemize}
  \item \textbf{Scored}: all six metrics are supported by the technique's sources
        (\LiveNissScored{} live techniques; these were scored before the rule was
        introduced and have not yet been re-examined against it).
  \item \textbf{Provisional}: some metrics are supported; the record stores those
        values, a source for each, and the lowest and highest score the formula could
        return (\LiveNissProvisional{}). A range of 0.4 to 9.4 says that one metric is
        known and five are not; it is a bound, not an estimate, and not a severity.
  \item \textbf{Unscored}: no peer-reviewed source, or no metric supported
        (\LiveNissUnscored{}).
\end{itemize}

\noindent Read the scoring state before the score. Three limits follow from the formula and
the data:

\begin{enumerate}
  \item \textbf{Impact only.} NISS has no likelihood term, so a technique with several
        disclosed vulnerabilities can score low.
  \item \textbf{Read-only harm cannot exceed 3.0.} With CR at Critical, CV at Implicit
        and nothing physical, the score is $(0.5 \cdot 10 + 10)/5 = 3.0$. NISS therefore
        rates every pure privacy-and-consent harm as low, whatever its scale. This is a
        property of the default weights, and it is the clearest reason to report the
        vector, not only the number.
  \item \textbf{No calibration or inter-rater study.} The levels were assigned by one
        author with AI assistance (Section~\ref{sec:limitations}).
\end{enumerate}

\noindent NISS scores signal-level interference with a person's neural function and the
consent under which it occurs; it does not measure mental states, and a score is not a
clinical or legal finding. Whether NISS should be reported as a single number or as its
vector with the PINS flag is an open question.

\subsection{NISS Severity Distribution}

The NISS severity distribution in the snapshot, as published in v1.5:

\begin{table}[H]
\centering
\caption{NISS severity distribution (\SnapshotTechniques{} snapshot techniques, as published in v1.5).}
\label{tab:niss-severity}
\small
\begin{tabular}{@{}l r r@{}}
\toprule
\textbf{NISS Severity} & \textbf{Count} & \textbf{\%} \\
\midrule
High     & 19 & 17.4\% \\
Medium   & 37 & 33.9\% \\
Low      & 52 & 47.7\% \\
None     &  1 &  0.9\% \\
\bottomrule
\end{tabular}
\end{table}

The registry file committed with v1.5 gives 21 high, 37 medium, 50 low and 1 none, and
several of its stored scores do not reproduce from Equation~\ref{eq:niss}; the live scores
do. The NISS distribution sits lower than the filed severity (Table~\ref{tab:severity})
because NISS separates purely digital impact from impact on a person.


\subsection{Neurological Extension (v1.1)}

NISS v1.0 mapped technique outcomes exclusively to DSM-5-TR psychiatric
categories. This created a systematic blind spot: neurological disruptions---loss
of smell, tinnitus, vestibular dysfunction, motor tremor, autonomic
dysregulation---have no DSM-5-TR entry because the DSM is a psychiatric manual,
not a neurological reference.

NISS v1.1 extends the impact chain to 42~neurological conditions drawn from
ICD-10-CM chapters G~(diseases of the nervous system), H~(diseases of the eye
and ear), and R~(symptoms and signs not elsewhere classified), using Adams \&
Victor's \textit{Principles of Neurology} (12th~ed.)~as the primary clinical
reference. These conditions are organized into seven categories:

\begin{itemize}
  \item \textbf{Sensory} (14 conditions): auditory (tinnitus, hearing loss),
    visual (scotoma, cortical blindness), vestibular (vertigo, nystagmus),
    olfactory (anosmia, parosmia), gustatory (ageusia, dysgeusia), and
    somatosensory (paresthesia, allodynia)
  \item \textbf{Motor} (8 conditions): tremor, dystonia, ataxia, spasticity,
    fasciculations, myoclonus, bradykinesia, dyskinesia
  \item \textbf{Autonomic} (6 conditions): orthostatic hypotension,
    hyperhidrosis, cardiac arrhythmia, thermoregulatory dysfunction,
    gastroparesis, neurogenic bladder
  \item \textbf{Consciousness} (4 conditions): syncope, absence seizures,
    altered sensorium, coma
  \item \textbf{Speech/Language} (4 conditions): expressive aphasia, receptive
    aphasia, dysarthria, alexia
  \item \textbf{Reflex} (3 conditions): hyperreflexia, areflexia, abnormal
    plantar response
  \item \textbf{Pain} (3 conditions): neuropathic pain, trigeminal neuralgia,
    central pain syndrome
\end{itemize}

Each condition maps to specific QIF hourglass bands, neural pathway
dependencies, and NISS metrics. The CD metric was broadened from ``Cognitive
Disruption'' to ``Cognitive/Functional Disruption'' to encompass sensory, motor,
and autonomic disruption alongside cognitive interference. This definitional
change preserves backward compatibility: all existing CD scores remain valid,
and the scoring formula is unchanged. The extension adds coverage for
neurological outcomes that were previously invisible to the scoring system.

The neurological conditions are queryable via KQL as the
\texttt{neurological\_conditions} table, enabling cross-correlation with TARA
techniques, neural pathways, and hourglass bands.
